% TODO checklist:
% - [x] Reviewed all security modules for threat coverage
% - [x] Inspected src/security/intent_filter.py (Attack pattern detection)
% - [x] Inspected src/security/output_guard.py (Query injection prevention)
% - [x] Reviewed Q&A.md (Threat scenarios and mitigations)
% - [x] Reviewed README.md (Security architecture overview)
% - [x] Considered multi-tenant isolation requirements

\section{Threat Model and Attack Scenarios}
\label{sec:threat-model}

The Cineca Agentic Platform operates in an enterprise environment where multiple threat actors may attempt to compromise the system. This section presents a comprehensive threat model that identifies potential attackers, their capabilities, attack vectors, and the platform's defensive measures.

\subsection{Attacker Profiles}
\label{subsec:attacker-profiles}

The platform considers three primary attacker profiles:

\subsubsection{Malicious Tenant}

A legitimate tenant who attempts to abuse their access or breach isolation boundaries.

\paragraph{Capabilities:}
\begin{itemize}
    \item Valid authentication credentials for their tenant
    \item Knowledge of API endpoints and data structures
    \item Ability to craft sophisticated queries and tool invocations
    \item Access to LLM-generated content within their tenant
\end{itemize}

\paragraph{Goals:}
\begin{itemize}
    \item Access other tenants' data
    \item Exhaust shared resources (denial of service)
    \item Extract sensitive information from the knowledge graph
    \item Escalate privileges within the platform
\end{itemize}

\subsubsection{Compromised User}

A legitimate user account that has been taken over by an attacker.

\paragraph{Capabilities:}
\begin{itemize}
    \item Valid session token for the compromised account
    \item All permissions assigned to the compromised user
    \item Access to session history and context
    \item Ability to invoke tools within permission scope
\end{itemize}

\paragraph{Goals:}
\begin{itemize}
    \item Exfiltrate data accessible to the user
    \item Perform destructive operations within user's scope
    \item Pivot to higher-privilege accounts
    \item Establish persistence for future access
\end{itemize}

\subsubsection{External Attacker}

An unauthenticated attacker probing the system from outside the trust boundary.

\paragraph{Capabilities:}
\begin{itemize}
    \item Network access to public API endpoints
    \item Knowledge of common API vulnerabilities
    \item Automated scanning and fuzzing tools
    \item No valid credentials initially
\end{itemize}

\paragraph{Goals:}
\begin{itemize}
    \item Bypass authentication to gain access
    \item Exploit API vulnerabilities for code execution
    \item Denial of service attacks
    \item Information disclosure from error messages
\end{itemize}

\subsection{Threat Categories and Attack Scenarios}
\label{subsec:threat-categories}

\subsubsection{Data Exfiltration}

\paragraph{Threat:} Unauthorized extraction of sensitive data from the knowledge graph or other data stores.

\paragraph{Attack Scenarios:}
\begin{itemize}
    \item \textbf{Cross-Tenant Query}: Crafting Cypher queries that bypass tenant filters to access other tenants' nodes and relationships.
    \item \textbf{Bulk Export}: Using legitimate graph tools to extract large volumes of data incrementally.
    \item \textbf{LLM Prompt Injection}: Manipulating prompts to coerce the LLM into revealing training data or system information.
\end{itemize}

\paragraph{Mitigations:}
\begin{itemize}
    \item Mandatory tenant ID filtering in all graph queries (enforced at the query generation layer)
    \item Result size limits enforced by the output guard (\texttt{LIMIT 100} default)
    \item Rate limiting to prevent high-volume extraction
    \item Intent filter detection of bulk export patterns (``dump all'', ``export everything'')
    \item Audit logging of all data access operations
\end{itemize}

\subsubsection{Privilege Escalation}

\paragraph{Threat:} Gaining access beyond the permissions granted to the authenticated user.

\paragraph{Attack Scenarios:}
\begin{itemize}
    \item \textbf{Role Confusion}: Exploiting JWT claim handling to inject admin roles.
    \item \textbf{Scope Injection}: Manipulating token payloads to include additional scopes.
    \item \textbf{Tool Permission Bypass}: Invoking tools without proper authorization checks.
\end{itemize}

\paragraph{Mitigations:}
\begin{itemize}
    \item Cryptographic JWT signature verification using JWKS public keys
    \item Server-side role-to-scope expansion (never trusting client-provided scopes directly)
    \item Per-tool authorization checks before invocation
    \item Audit logging of all authorization decisions
    \item Clear separation between role assignment (IdP) and scope expansion (platform)
\end{itemize}

\subsubsection{NL-to-Cypher Injection}

\paragraph{Threat:} Manipulating natural language inputs to generate malicious Cypher queries.

\paragraph{Attack Scenarios:}
\begin{itemize}
    \item \textbf{Destructive Queries}: Prompts designed to generate DELETE, DROP GRAPH, or TRUNCATE statements.
    \item \textbf{Schema Discovery}: Queries to expose the full graph schema for planning further attacks.
    \item \textbf{Unbounded Traversals}: Prompts that generate resource-exhausting path queries.
\end{itemize}

\paragraph{Mitigations:}
\begin{itemize}
    \item \textbf{Layer 1: Syntax Validation}: All generated Cypher is parsed before execution.
    \item \textbf{Layer 2: Read-Only Enforcement}: Write verbs (CREATE, MERGE, SET, DELETE, REMOVE) are blocked unless explicitly allowed.
    \item \textbf{Layer 3: Tenant Boundary Injection}: Tenant filters are automatically inserted into queries.
    \item \textbf{Layer 4: Depth Limits}: Path traversal depth is bounded (\texttt{-[*1..5]->} instead of \texttt{-[*]->}).
    \item \textbf{Layer 5: Timeout Guards}: Query execution has strict time limits.
    \item \textbf{Layer 6: Result Size Caps}: Maximum nodes and relationships per result set.
\end{itemize}

\subsubsection{Tool Abuse}

\paragraph{Threat:} Misusing MCP tools to perform unauthorized operations.

\paragraph{Attack Scenarios:}
\begin{itemize}
    \item \textbf{System Tool Exploitation}: Using system.backup or system.metrics to extract configuration.
    \item \textbf{Graph CRUD Abuse}: Bulk creating or deleting nodes to corrupt the knowledge graph.
    \item \textbf{Rate Limit Evasion}: Distributing requests across tools to bypass per-tool limits.
\end{itemize}

\paragraph{Mitigations:}
\begin{itemize}
    \item Per-tool required scopes enforced before invocation
    \item Role-based tool allow/deny lists in \texttt{roles.yaml}
    \item Per-tool rate limits (\texttt{tool\_rpm\_overrides})
    \item Input validation via JSON Schema for all tool parameters
    \item Audit logging of all tool invocations with input/output hashes
\end{itemize}

\subsubsection{Budget Abuse}

\paragraph{Threat:} Exploiting LLM capabilities to generate excessive API costs.

\paragraph{Attack Scenarios:}
\begin{itemize}
    \item \textbf{Token Exhaustion}: Submitting extremely long prompts to maximize token consumption.
    \item \textbf{Expensive Model Selection}: Choosing high-cost models for simple tasks.
    \item \textbf{Parallel Request Floods}: Overwhelming the system with concurrent LLM calls.
\end{itemize}

\paragraph{Mitigations:}
\begin{itemize}
    \item Per-tenant token limits (\texttt{tpm} in role configuration)
    \item Per-request maximum tokens (\texttt{max\_output\_tokens})
    \item Concurrent job limits (\texttt{concurrent\_jobs})
    \item Cost tracking and alerting per tenant
    \item Model access restrictions by role
\end{itemize}

\subsubsection{Session and State Attacks}

\paragraph{Threat:} Exploiting session management to gain unauthorized access or corrupt state.

\paragraph{Attack Scenarios:}
\begin{itemize}
    \item \textbf{Session Hijacking}: Stealing or replaying session tokens.
    \item \textbf{Session Fixation}: Forcing a known session ID on a victim.
    \item \textbf{Context Poisoning}: Injecting malicious content into session history.
\end{itemize}

\paragraph{Mitigations:}
\begin{itemize}
    \item Short JWT expiration times (configurable, enforced for internal endpoints)
    \item Cryptographic token signing (RS256/ES256 via JWKS)
    \item Session binding to tenant and user identity
    \item Session history stored server-side (Redis) not in client tokens
    \item PII scrubbing applied to session context before storage
\end{itemize}

\subsection{Defense Mapping}
\label{subsec:defense-mapping}

Table~\ref{tab:threat-defense-mapping} summarizes the mapping between threats and defensive controls.

\begin{table}[ht]
\centering
\caption{Threat-to-Defense Mapping}
\label{tab:threat-defense-mapping}
\small
\begin{tabular}{lp{7cm}}
\hline
\textbf{Threat} & \textbf{Primary Defenses} \\
\hline
Data Exfiltration & Tenant isolation, rate limits, result caps, audit logging \\
Privilege Escalation & JWT validation, RBAC, per-tool auth, scope expansion \\
NL-to-Cypher Injection & 6-layer validation pipeline, read-only mode, timeout guards \\
Tool Abuse & Required scopes, role-based allow/deny, per-tool rate limits \\
Budget Abuse & Token limits, cost tracking, concurrent job limits \\
Session Attacks & Short TTL, cryptographic signing, server-side state \\
Cross-Tenant Access & Mandatory tenant filters, key namespacing, allowlists \\
DoS/Resource Exhaustion & Rate limiting, circuit breakers, health probes \\
\hline
\end{tabular}
\end{table}

\subsection{Trust Boundaries}
\label{subsec:trust-boundaries}

The platform defines clear trust boundaries:

\begin{enumerate}
    \item \textbf{External Network $\rightarrow$ API Gateway}: All external traffic passes through authentication.
    \item \textbf{API Layer $\rightarrow$ Service Layer}: Authorization is enforced before service access.
    \item \textbf{Service Layer $\rightarrow$ Data Layer}: Tenant isolation is enforced at query time.
    \item \textbf{Platform $\rightarrow$ External LLM}: Circuit breakers and timeouts protect against provider issues.
    \item \textbf{Platform $\rightarrow$ Identity Provider}: JWKS validation ensures token integrity.
\end{enumerate}

\subsection{Residual Risks}
\label{subsec:residual-risks}

Despite comprehensive defenses, some residual risks remain:

\begin{itemize}
    \item \textbf{Zero-Day Vulnerabilities}: Unknown vulnerabilities in dependencies (FastAPI, python-jose, etc.).
    \item \textbf{LLM Unpredictability}: LLM outputs may occasionally bypass safety filters.
    \item \textbf{Insider Threats}: Operators with infrastructure access could bypass application-level controls.
    \item \textbf{Sophisticated Prompt Attacks}: Novel prompt injection techniques not yet recognized by filters.
\end{itemize}

\paragraph{Risk Acceptance:} These residual risks are managed through:
\begin{itemize}
    \item Regular dependency updates and security scanning
    \item Defense-in-depth with multiple overlapping controls
    \item Comprehensive audit logging for forensic analysis
    \item Incident response procedures for rapid mitigation
\end{itemize}

